% 書衣（後）：同色素面，訂線在左；下方小字為本版說明。
\begin{tikzpicture}[remember picture,overlay,shift={(current page.south west)},xshift=\edxshift,x=1mm,y=1mm]
  \node[anchor=south west,inner sep=0pt] at (0,0)
    {\includegraphics[width=160.9mm,height=254mm]{figures/cover_indigo.jpg}};
  \sewing{11}{-1}
  \node[text=coverpale,align=center,text width=110mm,inner sep=0pt,anchor=south] at (82,22)
    {\songR\fontsize{8.5pt}{14pt}\selectfont
     \XeTeXlinebreaklocale "zh"\XeTeXlinebreakskip=0pt plus 1pt
     清李漁撰《閒情偶寄》第\edslip 冊（\eddesc）　重排本\par
     底本：北京大學圖書館藏康熙間翼聖堂刻本　參校：哈佛燕京圖書館藏本\par
     錄文、校對、排版：Claude（Anthropic）輔助，二〇二六年\par
     原書文字屬公有領域；本版排版、校記以 \mbox{CC BY 4.0} 釋出\par
     正文句讀據維基文庫（\mbox{CC BY-SA 4.0}）\par
     \mbox{github.com/metachow/xianqing-ouji}\par};   % \mbox: no CJK break glue after the slashes
\end{tikzpicture}
